\section{Limitations}
\label{sec:conclusions_limitations}

Two main limitations qualify the scope of the findings presented in this thesis. Limitations specific to each contribution are already discussed in the corresponding chapters.

First, generalizability remains unverified. In principle, the approaches developed here are transferable but the analyses and models have not been validated in other cities or \acp{bss}, so empirical transferability to different urban contexts has not been assessed. Furthermore, the work is framed in dense, urban settings where \acp{bss} are typically deployed. Applicability to rural or low-density environments--where demand patterns, station spacing, and operational constraints differ substantially--has not been explored and remains an open question.

Second, all methods were evaluated on historical data only. The station planning framework has not been validated with stakeholder feedback, and the battery and maintenance predictions have not been tested in production. Assessing their real-world impact would require deployment and evaluation in live operator environments.
